\documentclass[10pt,a4paper]{amsart}
\usepackage{amsmath,amssymb,amsthm,graphicx}
\usepackage{verbatim}
\newenvironment{reference}{\comment}{\endcomment}
%\newenvironment{reference}{}{}
\newenvironment{slogan}{\comment}{\endcomment}
\newenvironment{history}{\comment}{\endcomment}
% For commutative diagrams we use Xy-pic
\usepackage[all]{xy}

% We use 2cell for 2-commutative diagrams.
\xyoption{2cell}
% Theorem environments.
%
\theoremstyle{plain}
\newtheorem{theorem}[subsection]{Theorem}
\newtheorem{proposition}[subsection]{Proposition}
\newtheorem{lemma}[subsection]{Lemma}

\theoremstyle{definition}
\newtheorem{definition}[subsection]{Definition}
\newtheorem{example}[subsection]{Example}
\newtheorem{exercise}[subsection]{Exercise}
\newtheorem{situation}[subsection]{Situation}

\theoremstyle{remark}
\newtheorem{remark}[subsection]{Remark}
\newtheorem{remarks}[subsection]{Remarks}

\numberwithin{equation}{subsection}

% Macros
%
\def\lim{\mathop{\mathrm{lim}}\nolimits}
\def\colim{\mathop{\mathrm{colim}}\nolimits}
\def\Spec{\mathop{\mathrm{Spec}}}
\def\Hom{\mathop{\mathrm{Hom}}\nolimits}
\def\Ext{\mathop{\mathrm{Ext}}\nolimits}
\def\SheafHom{\mathop{\mathcal{H}\!\mathit{om}}\nolimits}
\def\SheafExt{\mathop{\mathcal{E}\!\mathit{xt}}\nolimits}
\def\Sch{\mathit{Sch}}
\def\Mor{\mathop{\mathrm{Mor}}\nolimits}
\def\Ob{\mathop{\mathrm{Ob}}\nolimits}
\def\Sh{\mathop{\mathit{Sh}}\nolimits}
\def\NL{\mathop{N\!L}\nolimits}
\def\CH{\mathop{\mathrm{CH}}\nolimits}
\def\proetale{{pro\text{-}\acute{e}tale}}
\def\etale{{\acute{e}tale}}
\def\QCoh{\mathit{QCoh}}
\def\Ker{\mathop{\mathrm{Ker}}}
\def\Im{\mathop{\mathrm{Im}}}
\def\Coker{\mathop{\mathrm{Coker}}}
\def\Coim{\mathop{\mathrm{Coim}}}

% Boxtimes
%
\DeclareMathSymbol{\boxtimes}{\mathbin}{AMSa}{"02}

%
% Macros for moduli stacks/spaces
%
\def\QCohstack{\mathcal{QC}\!\mathit{oh}}
\def\Cohstack{\mathcal{C}\!\mathit{oh}}
\def\Spacesstack{\mathcal{S}\!\mathit{paces}}
\def\Quotfunctor{\mathrm{Quot}}
\def\Hilbfunctor{\mathrm{Hilb}}
\def\Curvesstack{\mathcal{C}\!\mathit{urves}}
\def\Polarizedstack{\mathcal{P}\!\mathit{olarized}}
\def\Complexesstack{\mathcal{C}\!\mathit{omplexes}}
% \Pic is the operator that assigns to X its picard group, usage \Pic(X)
% \Picardstack_{X/B} denotes the Picard stack of X over B
% \Picardfunctor_{X/B} denotes the Picard functor of X over B
\def\Pic{\mathop{\mathrm{Pic}}\nolimits}
\def\Picardstack{\mathcal{P}\!\mathit{ic}}
\def\Picardfunctor{\mathrm{Pic}}
\def\Deformationcategory{\mathcal{D}\!\mathit{ef}}
\usepackage{cleveref}
\crefname{originaldefinition}{original definition}{original definitions}
\Crefname{originaldefinition}{Original definition}{Original definitions}
\makeatletter
\newlabel{definition-fill-in-diagram}{{definition-fill-in-diagram}{source}{}{sourcepointer-1}{}}
\newlabel{definition-fill-in-diagram@cref}{{[originaldefinition][0][]definition-fill-in-diagram}{[1][0][]source}{}{}{}}
\makeatother
\crefname{originallemma}{original lemma}{original lemmas}
\Crefname{originallemma}{Original lemma}{Original lemmas}
\makeatletter
\newlabel{lemma-cat-dotted-arrows-base-change}{{lemma-cat-dotted-arrows-base-change}{source}{}{sourcepointer-2}{}}
\newlabel{lemma-cat-dotted-arrows-base-change@cref}{{[originallemma][0][]lemma-cat-dotted-arrows-base-change}{[1][0][]source}{}{}{}}
\makeatother
\crefname{originallemma}{original lemma}{original lemmas}
\Crefname{originallemma}{Original lemma}{Original lemmas}
\makeatletter
\newlabel{lemma-quasi-compact-permanence}{{lemma-quasi-compact-permanence}{source}{}{sourcepointer-3}{}}
\newlabel{lemma-quasi-compact-permanence@cref}{{[originallemma][0][]lemma-quasi-compact-permanence}{[1][0][]source}{}{}{}}
\makeatother
\crefname{originallemma}{original lemma}{original lemmas}
\Crefname{originallemma}{Original lemma}{Original lemmas}
\makeatletter
\newlabel{lemma-reach-points-scheme-theoretic-image}{{lemma-reach-points-scheme-theoretic-image}{source}{}{sourcepointer-4}{}}
\newlabel{lemma-reach-points-scheme-theoretic-image@cref}{{[originallemma][0][]lemma-reach-points-scheme-theoretic-image}{[1][0][]source}{}{}{}}
\makeatother
\title{Existence for the valuative criterion implies that the morphism is universally closed}
\author{The Stacks Project Authors}
\date{stacks-project:stacks-morphisms.tex:local-snapshot-20261012; GNU FDL 1.2 or later}
\begin{document}
\maketitle
\noindent\textit{Editorial note.} All delivered mathematical text is the author's own, copied from the pinned Stacks Project chapter source sources/stacks-algebraic-stacks-morphisms/pinned-original-source.tex; no statement or proof step was rewritten, reordered or removed.\par\medskip
\noindent\textit{Editorial note.} Preamble provenance: the Stacks Project keeps one single shared preamble for every chapter; this card takes it directly from the staged sources/stacks-algebraic-stacks-morphisms/preamble.tex as three verbatim spans (lines 8-12, 14-18 and 154-217), so no derived file was written and nothing was added or redefined. The wrapper class and the base packages are supplied by the assembler.\par\medskip
\noindent\textit{Editorial note.} Bibliography provenance: neither the counted statement nor its proof contains a citation (this lane rescans the whole span before writing the card), so the card carries no bibliography and no entry was reconstructed.\par\medskip
\noindent\textit{Editorial note.} Selection provenance: this card was opened under the file-based control contract (CONTROL.md), with the mechanical pre-card gate run on its statement and proof interval before the card was written.\par\medskip
\noindent\textit{Allowed standard prerequisites (editorial).}\begin{itemize}
\item The base change reduction for the category of dotted arrows. Invoked at the start of the counted proof (source lines 9794-9797). Allows the diagram to be replaced by one over a suitable base change.
\item The description of the category of dotted arrows and of its set of isomorphism classes. Used throughout the counted proof (source lines 9798-9820). Turns the existence hypothesis into the existence of a dotted arrow.
\item The statement that a morphism satisfying existence in the valuative criterion is universally closed. Invoked near the end of the counted proof (source lines 9821-9832). Supplies the conclusion once the hypothesis has been reformulated.
\end{itemize}\medskip
\ifdefined\hypertarget\hypertarget{sourcepointer-1}{}\fi
\noindent\textit{Source reference.} Original definition definition-fill-in-diagram: sources/stacks-algebraic-stacks-morphisms/pinned-original-source.tex, source lines 9061--9061. Navigation pointer only: the counted text refers to the Stacks item definition-fill-in-diagram, which is an existing labeled result of this same chapter and is not part of this card. The pointer carries the target's own label as its number and its definition line as its span; its content is not reproduced here.\par\medskip
\ifdefined\hypertarget\hypertarget{sourcepointer-2}{}\fi
\noindent\textit{Source reference.} Original lemma lemma-cat-dotted-arrows-base-change: sources/stacks-algebraic-stacks-morphisms/pinned-original-source.tex, source lines 9187--9187. Navigation pointer only: the counted text refers to the Stacks item lemma-cat-dotted-arrows-base-change, which is an existing labeled result of this same chapter and is not part of this card. The pointer carries the target's own label as its number and its definition line as its span; its content is not reproduced here.\par\medskip
\ifdefined\hypertarget\hypertarget{sourcepointer-3}{}\fi
\noindent\textit{Source reference.} Original lemma lemma-quasi-compact-permanence: sources/stacks-algebraic-stacks-morphisms/pinned-original-source.tex, source lines 2044--2044. Navigation pointer only: the counted text refers to the Stacks item lemma-quasi-compact-permanence, which is an existing labeled result of this same chapter and is not part of this card. The pointer carries the target's own label as its number and its definition line as its span; its content is not reproduced here.\par\medskip
\ifdefined\hypertarget\hypertarget{sourcepointer-4}{}\fi
\noindent\textit{Source reference.} Original lemma lemma-reach-points-scheme-theoretic-image: sources/stacks-algebraic-stacks-morphisms/pinned-original-source.tex, source lines 2101--2101. Navigation pointer only: the counted text refers to the Stacks item lemma-reach-points-scheme-theoretic-image, which is an existing labeled result of this same chapter and is not part of this card. The pointer carries the target's own label as its number and its definition line as its span; its content is not reproduced here.\par\medskip
\begin{lemma}
\label{lemma-converse-existence-universally-closed}
Let $f : \mathcal{X} \to \mathcal{Y}$ be a morphism of algebraic stacks.
Assume
\begin{enumerate}
\item $f$ is quasi-separated, and
\item $f$ is universally closed.
\end{enumerate}
Then $f$ satisfies the existence part of the valuative criterion.
\end{lemma}

\begin{proof}
Consider a solid diagram
$$
\xymatrix{
\Spec(K) \ar[r]_-x \ar[d]_j & \mathcal{X} \ar[d]^f \\
\Spec(A) \ar[r]^-y \ar@{..>}[ru] & \mathcal{Y}
}
$$
where $A$ is a valuation ring with field of fractions $K$
and $\gamma : y \circ j \longrightarrow f \circ x$ as in
Definition \ref{definition-fill-in-diagram}. By
Lemma \ref{lemma-cat-dotted-arrows-base-change}
in order to find a dotted arrow (after possibly replacing
$K$ by an extension and $A$ by a valuation ring dominating it)
we may replace $\mathcal{Y}$ by $\Spec(A)$ and $\mathcal{X}$
by $\Spec(A) \times_\mathcal{Y} \mathcal{X}$. Of course
we use here that being
quasi-separated and universally closed are preserved under base change.
Thus we reduce to the case discussed in the next paragraph.

\medskip\noindent
Consider a solid diagram
$$
\xymatrix{
\Spec(K) \ar[r]_-x \ar[d]_j & \mathcal{X} \ar[d]^f \\
\Spec(A) \ar@{=}[r] \ar@{..>}[ru] & \Spec(A)
}
$$
where $A$ is a valuation ring with field of fractions $K$ as in
Definition \ref{definition-fill-in-diagram}.
By Lemma \ref{lemma-quasi-compact-permanence} and the fact
that $f$ is quasi-separated we have that
the morphism $x$ is quasi-compact.
Since $f$ is universally closed, we have in particular
that $|f|(\overline{\{x\}})$ is closed in $\Spec(A)$.
Since this image contains the generic point of $\Spec(A)$
there exists a point $x' \in |\mathcal{X}|$ in the closure
of $x$ mapping to the closed point of $\Spec(A)$.
By Lemma \ref{lemma-reach-points-scheme-theoretic-image}
we can find a commutative diagram
$$
\xymatrix{
\Spec(K') \ar[r] \ar[d] & \Spec(K) \ar[d] \\
\Spec(A') \ar[r] & \mathcal{X}
}
$$
such that the closed point of $\Spec(A')$ maps to $x' \in |\mathcal{X}|$.
It follows that $\Spec(A') \to \Spec(A)$ maps the closed point
to the closed point, i.e., $A'$ dominates $A$ and this finishes the proof.
\end{proof}

\end{document}
